% Frozen P55 protocol described for the provisional manuscript.
\subsection{一覧と詳細を分けた表示比較}
API説明の省略による変化を調べるため，全表示F，詳細のみ省略D，
同じ詳細と一覧を省略Bを設けた．
古い成功出力のうち，保存したAPI返却と表示処理から出自を照合できる区間だけを変更する．
操作コード，リテラル，混合出力，エラー，最新応答を保持し，
DとBで詳細の省略集合と省略印の形式をそろえる．
Python状態，公開スキーマ，原文の参照経路，公開契約の拒否時返信は共通である．
従って，実行環境から知識を消去する介入ではない．

以前に観測済みの４課題の第８応答後を使い，各３条件を一度ずつ実行する計画を固定した．
上限は各16生成，継続120,000 token，480秒とし，初回は入力キャッシュを再利用しない．
全入力と残予算を計数して出力枠を制限する．
独立評価器で保存環境の最終成否を判定し，仕事全体の論理トークンを
$T_{\rm prefix}+\sum_j(I_j+O_j)$とする．
制御，計数，失敗，再取得を含め，復元・評価・推論器常駐は別記する．
使用量不明の生成は実行も再試行もせず当該条件を終了する．
API作用と片付けが既知の場合のみ最終状態を評価するが，不明費用は補完しない．

推論にはMiniMax-M3のIQ4\_XS量子化版とllama.cpp b11669を用いる．
コンテキスト上限は65,536，温度1.0，top-p 0.95，出力上限4,096とした．
継続のseedは50から呼出しごとに１増やし，条件間でそろえる．
API上限には前半の要求も含め，写真は160，他の課題は80とする．
順序は課題ごとにFDB，DBF，BFD，FBDで固定した．
量子化版，推論器版，実装版，入力と返却のハッシュを保存し，
実行条件を固定したことと結果の決定性を区別する．

主要な対照はD対FとD対Bである．
返金課題には追加で省ける一覧がなく，D/Bの基礎入力が同じ負の対照となる．
ここで生じた差は一覧保持の効果と解釈しない．
既知４課題での開発比較のため，独立12課題や統計的な非劣性の証明とは扱わない．
